% Cote 119: offprints, reports, courses, CV and letters (typed and printed part): archive cote 119, site page 237 (PDF page 238).
% typescript. Draft compared: none.
% AI-transcribed by Claude Opus 5.5 (claude-opus-5-5, adaptive thinking, effort high), Message Batches API, batch r119_full (msgbatch_01DwZNm5Y7YsKBWyAN4sFs8w), 2026-11-13; not yet checked by eye.
\scanpage{119/237}
\mspage{- 25 -}

LISTE DES TRAVAUX CITES, SUSCITES OU INFLUENCES PAR
LES TRAVAUX DE A. GROTHENDIECK.

M. ARTIN, Algebraization of Formal Moduli, I (in
Global Analysis, pp. 21-71, University
of Tokyo Press, Princeton University Press,
1968), II Existence of modifications,
Annals of Mathematics, Vol. 91, pp. 88-135
(1970).

M.F. ATIYAH et F. HIRZEBRUCH, Riemann-Roch theorems
for differentiable manifolds, Bull. Amer.
Math. Soc., vol. 65, pp. 276-281 (1959).

P. BERTHELOT, Cohomologie cristalline des schémas
propres et lisses de caractéristique $p > 0$,
Thèse, Université Paris VII, 1971 (paraîtra
dans Lecture Notes of Math. chez Springer).

P. BERTHELOT et L. ILLUSIE, Classes de Chern en
cohomologie cristalline, C.R. Acad. Sci.
Paris t. 270, pp. 1695-1697 (22 juin 1970)
et p. 1750-1752 (29 juin 1970).

A. BOREL et J.P. SERRE, Le théorème de Riemann-Roch
(d'après des résultats inédits de
A. Grothendieck), Bull. Soc. Math. France,
t. 86, pp. 97-136 (1958).

P. DELIGNE, Théorie de Hodge I (Actes du Congrès
International des mathématiciens, Nice 1970)
et II, Publications Math. n$^\circ$ 40, pp. 5-57
(1971).

P. DELIGNE et D. MUMFORD, The irreducibility of
the space of curves of a given genus,
Pub. Math. n$^\circ$ 36, pp. 75-110 (1969).

\margin{237}
